% Part 2 — Challenge Corner

\begin{problem}[The mathematics of Benford's law]
Many real-world data sets, such as populations, lengths, financial amounts, and physical constants, have leading digits that are not uniformly distributed. Benford's law models the leading digit $D\in\{1,\ldots,9\}$ by
\[
P(D=d)=\log_{10}\!\left(1+\frac1d\right).
\]
In this problem you will verify that this formula really defines a probability distribution and then connect it to a simple continuous model.
\begin{enumerate}[label=\textbf{(\roman*)}]
\item Show that $\sum_{d=1}^{9}P(D=d)=1$, so the formula is a valid probability mass function.
\item Show that the probability of a leading digit $1$, $2$, or $3$ is greater than $60\%$.
\item By expanding $\sum dP(D=d)$ and regrouping logarithmic terms, prove that the mean leading digit is $\mu=9-\log_{10}(9!)$.
\item If $U\sim \mathrm{Unif}[0,1)$ and $Y=10^{U}$, show that the leading digit of $Y$ follows Benford's law.
\end{enumerate}
\end{problem}

\begin{hint}
\begin{itemize}[leftmargin=*]
\item $\log(1+\tfrac1d)=\log(d+1)-\log d$.
\item Sum three explicit logs to compare to $\log 4$.
\item Expand expectations and regroup coefficients of $\log k$.
\item Leading digit $d$ iff $d\le Y<d+1$; take $\log_{10}$.
\end{itemize}
\end{hint}

\begin{solution}
\textbf{(i)} Rewrite each probability as
\[
\log_{10}\!\left(1+\frac1d\right)
=\log_{10}\!\left(\frac{d+1}{d}\right)
=\log_{10}(d+1)-\log_{10}d.
\]
Therefore the sum telescopes:
\[
\sum_{d=1}^{9}P(D=d)
=\log_{10}2-\log_{10}1+\cdots+\log_{10}10-\log_{10}9
=\log_{10}10-\log_{10}1=1.
\]

\textbf{(ii)} The same telescoping idea gives
\[
P(D\in\{1,2,3\})=\log_{10}2+\log_{10}\frac32+\log_{10}\frac43
=\log_{10}4\approx0{.}602.
\]
Thus the first three digits account for slightly more than $60\%$ of all leading digits under Benford's law.

\textbf{(iii)} The mean is
\[
E(D)=\sum_{d=1}^{9}d\log_{10}\!\left(\frac{d+1}{d}\right)
=\sum_{d=1}^{9}d\bigl(\log_{10}(d+1)-\log_{10}d\bigr).
\]
Expanding shows that most logarithms appear twice with neighbouring coefficients:
\[
E(D)=1\log_{10}2+2\log_{10}3-2\log_{10}2+\cdots+9\log_{10}10-9\log_{10}9.
\]
After regrouping, this becomes
\[
E(D)=9\log_{10}10-\log_{10}(1\cdot2\cdots9)
=9-\log_{10}(9!).
\]

\textbf{(iv)} The leading digit of $Y$ is $d$ precisely when
\[
d\le Y<d+1.
\]
Since $Y=10^U$, taking $\log_{10}$ gives
\[
\log_{10}d\le U<\log_{10}(d+1).
\]
Because $U$ is uniform on $[0,1)$, the probability is the length of this interval:
\[
P(D=d)=\log_{10}(d+1)-\log_{10}d
=\log_{10}\!\left(1+\frac1d\right).
\]
\end{solution}

\begin{takeaways}
\item Logarithm differences often telescope into compact probability sums.
\item Benford leading digits are heavily weighted toward smaller digits.
\item Uniformity on a logarithmic scale produces Benford's law on the original scale.
\end{takeaways}

\begin{problem}[The squared Cauchy distribution]
Let
\[
g(x)=\frac{1}{(1+x^2)^2},\qquad -\infty<x<\infty.
\]
This density-like curve has heavier tails than a normal distribution but still has finite mean and variance. The aim is to normalise it into a probability density function and then compute its first two moments carefully.
\begin{enumerate}[label=\textbf{(\roman*)}]
\item Show $\displaystyle\int_{-\infty}^{\infty}g(x)\,dx=\frac{\pi}{2}$ using the substitution $x=\tan\theta$.
\item Verify that $f(x)=\dfrac{2}{\pi(1+x^2)^2}$ is a valid probability density function on the real line.
\item Show $E(X)=0$, and justify why the improper integral converges before using symmetry.
\item Show that the variance is defined and that $\mathrm{Var}(X)=1$.
\end{enumerate}
\end{problem}

\begin{hint}
\begin{itemize}[leftmargin=*]
\item Half-angle formulas after angular substitution.
\item Non-negativity and total mass $(2/\pi)\cdot\pi/2=1$.
\item Split $[0,\infty)$, substitute $u=1+x^{2}$ for the positive half before symmetrising.
\item $E[X^{2}]$ via same substitution yields $1$, then subtract $0^{2}$.
\end{itemize}
\end{hint}

\begin{solution}
\textbf{(i)} Let $x=\tan\theta$. Then $dx=\sec^2\theta\,d\theta$ and $1+x^2=\sec^2\theta$. As $x$ runs from $-\infty$ to $\infty$, $\theta$ runs from $-\frac{\pi}{2}$ to $\frac{\pi}{2}$. Hence
\[
\int_{-\infty}^{\infty}\frac{1}{(1+x^2)^2}\,dx
=\int_{-\pi/2}^{\pi/2}\frac{\sec^2\theta}{\sec^4\theta}\,d\theta
=\int_{-\pi/2}^{\pi/2}\cos^2\theta\,d\theta
=\frac{\pi}{2}.
\]

\textbf{(ii)} The function is non-negative for all real $x$. Also,
\[
\int_{-\infty}^{\infty}f(x)\,dx
=\frac{2}{\pi}\int_{-\infty}^{\infty}\frac{1}{(1+x^2)^2}\,dx
=\frac{2}{\pi}\cdot\frac{\pi}{2}=1.
\]
So $f$ is a valid pdf.

\textbf{(iii)} The integrand for the mean is
\[
xf(x)=\frac{2x}{\pi(1+x^2)^2},
\]
which is odd. Before using symmetry, check convergence on the positive half-line:
\[
\int_0^\infty \frac{2x}{\pi(1+x^2)^2}\,dx
=\frac{1}{\pi}\int_1^\infty u^{-2}\,du
=\frac1\pi,
\]
where $u=1+x^2$. The corresponding negative-half integral is finite and equal to $-\frac1\pi$, so $E(X)=0$.

\textbf{(iv)} Since $E(X)=0$, the variance is $E(X^2)$. Compute
\[
E(X^2)=\int_{-\infty}^{\infty}\frac{2x^2}{\pi(1+x^2)^2}\,dx.
\]
Using $x=\tan\theta$ again,
\[
\frac{2}{\pi}\int_{-\pi/2}^{\pi/2}\tan^2\theta\cos^2\theta\,d\theta
=\frac{2}{\pi}\int_{-\pi/2}^{\pi/2}\sin^2\theta\,d\theta
=1.
\]
Therefore $\mathrm{Var}(X)=1-0^2=\boxed{1}$.
\end{solution}

\begin{takeaways}
\item A normalising constant is chosen so the total area under a pdf is exactly $1$.
\item Symmetry is powerful, but improper integrals still need convergence checks.
\item The substitution $x=\tan\theta$ is natural whenever $1+x^2$ appears.
\end{takeaways}
